\section{Convert one validated triplet at a time}
The decoder keeps metadata in a private stack object while it calls the parser. Retain
the framebuffer pointer and capacity in saved registers: parser temporaries are not
part of the calling contract. After parsing succeeds, compare the pixel count with
unsigned capacity. A set high bit must not make a large capacity appear too small.

At the first pixel, the source cursor is input base plus raster offset. Its destination
cursor names the first output word. Keep a separate remaining-pixel count. Each
iteration reads three unsigned samples, forms the numerical color word and stores it,
then advances the cursors by three and four respectively. The validated dimensions
make the first remaining count positive. Decrement once per pixel and stop at zero.

The final cursors identify different object ends because their strides differ. Inspect
the last stored word and the following sentinel word. Correct colors alone do not
establish that the decoder stopped at the correct extent. Pack the width and height
for the return after conversion; neither the input byte count nor the framebuffer byte
count is a substitute for a dimension.

\section{Size the canonical header before storing it}
The canonical header has a fixed nine-byte contribution plus the decimal digit counts
of width and height. Each digit count is one or two under the exercise bounds. Thus
\[
 L=9+\operatorname{digits}(w)+\operatorname{digits}(h)+3wh.
\]
For the worked image, \(L=23\). At sixty-four by sixty-four it is 12301 bytes.
Validate dimensions first, then compute this bounded length and compare it with capacity
before any pixel read or header store. A short destination should not contain a new
magic number followed by an old raster.

Write the magic and LF, width, space, height, LF, then \texttt{255} and LF. For a
two-digit dimension, unsigned division by ten supplies the tens digit and remainder
supplies the units digit. Add ASCII zero to each numerical digit. A one-digit dimension
omits the zero quotient rather than writing a leading zero. Width nine and width ten
therefore exercise different header lengths even for the same height.

Once the header is complete, extract one input word's red, green and blue fields into
three consecutive bytes. The high byte does not enter the file. A byte-oriented encoder
can instead load the appropriate three byte positions from each little-endian word.
Both methods must preserve the sample codes and emit the same canonical bytes. The
completed-interface checks should accept either method.
